\documentclass[10pt,a4paper]{amsart}
\usepackage{amsmath,amssymb,amsthm,graphicx}
% Retained preamble for the Stacks Project chapter {CH} (shared project preamble).
% The chapter file pinned-original-source.tex only carries \input{{preamble}} on its first line,
% so no preamble text lives inside the chapter itself.  The Stacks Project uses one single shared
% preamble for every chapter: the snapshot stages it three times (sources/stacks-chow/preamble.tex,
% sources/stacks-etale-cohomology/preamble.tex and focused-seeds/AGP/sources/stacks-snapshot-20261011/
% preamble.tex) and those copies are identical in size.  The block below is copied mechanically,
% line for line, from sources/stacks-chow/preamble.tex lines 14-18 and 154-217: the Xy-pic setup and
% the native theorem declarations together with the mathematical macros.  Nothing was added, removed
% or redefined; the wrapper class and the amsmath/amssymb/amsthm/graphicx base packages are supplied
% by the assembler.
% For commutative diagrams we use Xy-pic
% Chapter: stacks-cohomology.tex -- Cohomology of Sheaves.
\usepackage[all]{xy}

% We use 2cell for 2-commutative diagrams.
\xyoption{2cell}

% Theorem environments.
%
\theoremstyle{plain}
\newtheorem{theorem}[subsection]{Theorem}
\newtheorem{proposition}[subsection]{Proposition}
\newtheorem{lemma}[subsection]{Lemma}

\theoremstyle{definition}
\newtheorem{definition}[subsection]{Definition}
\newtheorem{example}[subsection]{Example}
\newtheorem{exercise}[subsection]{Exercise}
\newtheorem{situation}[subsection]{Situation}

\theoremstyle{remark}
\newtheorem{remark}[subsection]{Remark}
\newtheorem{remarks}[subsection]{Remarks}

\numberwithin{equation}{subsection}

% Macros
%
\def\lim{\mathop{\mathrm{lim}}\nolimits}
\def\colim{\mathop{\mathrm{colim}}\nolimits}
\def\Spec{\mathop{\mathrm{Spec}}}
\def\Hom{\mathop{\mathrm{Hom}}\nolimits}
\def\Ext{\mathop{\mathrm{Ext}}\nolimits}
\def\SheafHom{\mathop{\mathcal{H}\!\mathit{om}}\nolimits}
\def\SheafExt{\mathop{\mathcal{E}\!\mathit{xt}}\nolimits}
\def\Sch{\mathit{Sch}}
\def\Mor{\mathop{\mathrm{Mor}}\nolimits}
\def\Ob{\mathop{\mathrm{Ob}}\nolimits}
\def\Sh{\mathop{\mathit{Sh}}\nolimits}
\def\NL{\mathop{N\!L}\nolimits}
\def\CH{\mathop{\mathrm{CH}}\nolimits}
\def\proetale{{pro\text{-}\acute{e}tale}}
\def\etale{{\acute{e}tale}}
\def\QCoh{\mathit{QCoh}}
\def\Ker{\mathop{\mathrm{Ker}}}
\def\Im{\mathop{\mathrm{Im}}}
\def\Coker{\mathop{\mathrm{Coker}}}
\def\Coim{\mathop{\mathrm{Coim}}}

% Boxtimes
%
\DeclareMathSymbol{\boxtimes}{\mathbin}{AMSa}{"02}

%
% Macros for moduli stacks/spaces
%
\def\QCohstack{\mathcal{QC}\!\mathit{oh}}
\def\Cohstack{\mathcal{C}\!\mathit{oh}}
\def\Spacesstack{\mathcal{S}\!\mathit{paces}}
\def\Quotfunctor{\mathrm{Quot}}
\def\Hilbfunctor{\mathrm{Hilb}}
\def\Curvesstack{\mathcal{C}\!\mathit{urves}}
\def\Polarizedstack{\mathcal{P}\!\mathit{olarized}}
\def\Complexesstack{\mathcal{C}\!\mathit{omplexes}}
% \Pic is the operator that assigns to X its picard group, usage \Pic(X)
% \Picardstack_{X/B} denotes the Picard stack of X over B
% \Picardfunctor_{X/B} denotes the Picard functor of X over B
\def\Pic{\mathop{\mathrm{Pic}}\nolimits}
\def\Picardstack{\mathcal{P}\!\mathit{ic}}
\def\Picardfunctor{\mathrm{Pic}}
\def\Deformationcategory{\mathcal{D}\!\mathit{ef}}
\usepackage{cleveref}
\crefname{originallemma}{original lemma}{original lemmas}
\Crefname{originallemma}{Original lemma}{Original lemmas}
\makeatletter
\newlabel{derived-lemma-derived-canonical-delta-functor}{{0152}{source}{}{sourcepointer-1}{}}
\newlabel{derived-lemma-derived-canonical-delta-functor@cref}{{[originallemma][0][]0152}{[1][0][]source}{}{}{}}
\makeatother
\crefname{originalsection}{original section}{original sections}
\Crefname{originalsection}{Original section}{Original sections}
\makeatletter
\newlabel{derived-section-canonical-delta-functor}{{014Z}{source}{}{sourcepointer-2}{}}
\newlabel{derived-section-canonical-delta-functor@cref}{{[originalsection][0][]014Z}{[1][0][]source}{}{}{}}
\makeatother
\crefname{originalsection}{original section}{original sections}
\Crefname{originalsection}{Original section}{Original sections}
\makeatletter
\newlabel{section-properties-K-injective}{{0D5U}{source}{}{sourcepointer-3}{}}
\newlabel{section-properties-K-injective@cref}{{[originalsection][0][]0D5U}{[1][0][]source}{}{}{}}
\makeatother
\crefname{originallemma}{original lemma}{original lemmas}
\Crefname{originallemma}{Original lemma}{Original lemmas}
\makeatletter
\newlabel{sheaves-lemma-j-shriek-modules}{{00A7}{source}{}{sourcepointer-4}{}}
\newlabel{sheaves-lemma-j-shriek-modules@cref}{{[originallemma][0][]00A7}{[1][0][]source}{}{}{}}
\makeatother
\title{Mayer-Vietoris distinguished triangle for extension by zero along the inclusions of two open subspaces}
\author{The Stacks Project Authors}
\date{stacks-project:cohomology.tex:local-snapshot-20261011; GNU FDL 1.2 or later}
\begin{document}
\maketitle
\noindent\textit{Editorial note.} All delivered mathematical text is the author's own, copied from the pinned Stacks Project chapter file cohomology.tex; no statement or proof step was rewritten, reordered or removed.\par\medskip
\noindent\textit{Editorial note.} The proof ends with the author's own '(computation omitted)' after naming the stalk computation; the parenthetical is reproduced verbatim. The omitted computation is a bounded stalk check which the named reference carries out, and the identification of the triangle with a short exact sequence is complete as written.\par\medskip
\noindent\textit{Editorial note.} The counted spans carry 4 navigation stubs and no bibliography citation.\par\medskip
\noindent\textit{Editorial note.} The preamble is the chapter's retained minimal preamble (sources/stacks-cohomology/retained-preamble.tex lines 1-81), a verbatim excerpt of the native preamble.\par\medskip
\noindent\textit{Editorial note.} The label lemma-exact-sequence-lower-shriek is defined only in cohomology.tex, so no cross-chapter restatement had to be weighed for this card.\par\medskip
\noindent\textit{Allowed standard prerequisites (editorial).}\begin{itemize}
\item Restriction and extension by zero along an open immersion are computed without deriving. Same chapter, Section section-properties-K-injective, invoked at proof lines 8098-8101. Supplies that an arbitrary complex represents the derived functors involved.
\item Stalk description of extension by zero for modules. Sheaves, Lemma sheaves-lemma-j-shriek-modules, invoked at proof lines 8113-8114. Supplies the termwise exactness check of the displayed sequence.
\item Distinguished triangle attached to a short exact sequence of complexes. Derived Categories, Section derived-section-canonical-delta-functor and Lemma derived-lemma-derived-canonical-delta-functor, invoked at proof lines 8102-8106. Supplies the triangle once the short exact sequence is in hand.
\end{itemize}\medskip
\ifdefined\hypertarget\hypertarget{sourcepointer-1}{}\fi
\noindent\textit{Source reference.} Original lemma 0152: sources/stacks-chow/tags.txt, source lines 1493--1493. Pointer outside the counted span, invoked once for context; the load-bearing argument of the counted statement is the author's proof reproduced in full in the proof span.\par\medskip
\ifdefined\hypertarget\hypertarget{sourcepointer-2}{}\fi
\noindent\textit{Source reference.} Original section 014Z: sources/stacks-chow/tags.txt, source lines 1490--1490. Pointer outside the counted span, invoked once for context; the load-bearing argument of the counted statement is the author's proof reproduced in full in the proof span.\par\medskip
\ifdefined\hypertarget\hypertarget{sourcepointer-3}{}\fi
\noindent\textit{Source reference.} Original section 0D5U: sources/stacks-cohomology/pinned-original-source.tex, source lines 7814--7814. Pointer outside the counted span, invoked once for context; the load-bearing argument of the counted statement is the author's proof reproduced in full in the proof span.\par\medskip
\ifdefined\hypertarget\hypertarget{sourcepointer-4}{}\fi
\noindent\textit{Source reference.} Original lemma 00A7: sources/stacks-chow/tags.txt, source lines 379--379. Pointer outside the counted span, invoked once for context; the load-bearing argument of the counted statement is the author's proof reproduced in full in the proof span.\par\medskip
\begin{lemma}
\label{lemma-exact-sequence-lower-shriek}
Let $(X, \mathcal{O}_X)$ be a ringed space.
Let $X = U \cup V$ be the union of two open subspaces.
For any object $E$ of $D(\mathcal{O}_X)$ we have a distinguished
triangle
$$
j_{U \cap V!}E|_{U \cap V} \to
j_{U!}E|_U \oplus j_{V!}E|_V \to E \to 
j_{U \cap V!}E|_{U \cap V}[1]
$$
in $D(\mathcal{O}_X)$.
\end{lemma}

\begin{proof}
We have seen in Section \ref{section-properties-K-injective}
that the restriction functors and the extension
by zero functors are computed by just applying the functors to
any complex. Let $\mathcal{E}^\bullet$ be a complex of $\mathcal{O}_X$-modules
representing $E$. The distinguished triangle of the lemma is the
distinguished triangle associated (by
Derived Categories, Section
\ref{derived-section-canonical-delta-functor} and especially
Lemma \ref{derived-lemma-derived-canonical-delta-functor})
to the short exact sequence of complexes of $\mathcal{O}_X$-modules
$$
0 \to j_{U \cap V!}\mathcal{E}^\bullet|_{U \cap V} \to
j_{U!}\mathcal{E}^\bullet|_U \oplus j_{V!}\mathcal{E}^\bullet|_V
\to \mathcal{E}^\bullet \to 0
$$
To see this sequence is exact one checks on stalks using
Sheaves, Lemma \ref{sheaves-lemma-j-shriek-modules}
(computation omitted).
\end{proof}

\end{document}
